
\section{Architecture Integration}

\subsection{Model Support}

\begin{table}[H]
\centering
\small
\begin{tabular}{llll}
\toprule
\textbf{Architecture} & \textbf{Nested?} & \textbf{RoPE?} & \textbf{GQA?} \\
\midrule
DistilBERT & No & No & No \\
RoBERTa & Yes & No & No \\
GPT-2 & No & No & No \\
Qwen3.5 & No & Yes & Yes \\
\bottomrule
\end{tabular}
\end{table}

\subsection{Weight Transfer}

Standard four-architecture copy with additional factorization step:

\begin{lstlisting}
with torch.no\_grad():
    for proj\_key, target\_mod in target\_map.items():
        w, b = extracted[proj\_key]
        if w is None:
            continue
        if w.ndim != 2:
            w = w.reshape(w.size(0), -1)
        if w.shape != (target\_mod.out\_features, target\_mod.in\_features) \
           and (target\_mod.out\_features, target\_mod.in\_features) == w.t().shape:
            w = w.t()
        target\_mod.decompose\_from\_weight(w[:target\_mod.out\_features, :target\_mod.in\_features], ...)
\end{lstlisting}

The \texttt{decompose\_from\_weight} method:
\begin{enumerate}
\item Computes SVD.
\item Selects rank by energy threshold.
\item Sets \texttt{down.weight = $\Sigma$^{1/2} V^T}, \texttt{up.weight = U $\Sigma$^{1/2}}.

\end{enumerate}
\subsection{GQA Handling}

For Qwen with $H_{kv} < H$, the K and V projections have smaller output 
dimensions. The factorization is done at the appropriate rank.

\begin{lstlisting}
self.k\_proj = FactorisedLinear(hidden\_size, self.all\_kv\_head\_size, ...)
self.v\_proj = FactorisedLinear(hidden\_size, self.all\_kv\_head\_size, ...)
\end{lstlisting}

\subsection{The `copy\_weights` Failure Path}

\textbf{Case 1: UniversalAdapter succeeds for all projections.}
\begin{itemize}
\item Standard extraction works.
\item No issue.

\end{itemize}
\textbf{Case 2: UniversalAdapter fails, but the direct extraction already succeeded.}
\begin{itemize}
\item \texttt{missing} is empty $\to$ no try/except triggered.
\item No issue.

\end{itemize}
\textbf{Case 3: UniversalAdapter raises.}
\begin{itemize}
\item \texttt{raise RuntimeError} aborts \texttt{copy\_weights}.
\item \texttt{inject\_orion} catches it (the outer try/except) $\to$ returns None.
\item \textbf{Safe} -- the module is not installed.

\end{itemize}
\textbf{Case 4: UniversalAdapter returns None for all \texttt{missing}.}
\begin{itemize}
\item \texttt{missing\_after} is non-empty $\to$ RuntimeError.
\item Same as case 3.

\end{itemize}
\textbf{Case 5: UniversalAdapter raises but injection proceeds anyway.}
\begin{itemize}
\item Not possible -- \texttt{inject\_orion} has its own try/except.

\end{itemize}
\subsection{Injection}

\textbf{Algorithm 2:} Inject ORION
\begin{lstlisting}
Input:  layer  , config
Output: new attention in  

1: src $\leftarrow$ extract\_attention( )
2: new $\leftarrow$ ORIONAttention(config, source\_attn=src)
3: if new.copy\_weights fails: exception $\to$ return None
4: install new in  .attention.self (or .attn, .self\_attn)
5: return new
\end{lstlisting}

\subsection{Common Pitfalls}

\subsubsection{`FactorisedLinear.weight` Materializes Full Matrix}

\begin{lstlisting}
@property
def weight(self):
    return torch.matmul(self.up.weight, self.down.weight)
\end{lstlisting}

For 768$\times$768 this is fine. For 3072$\times$768 with rank 64: \texttt{64 $\times$ 3072 + 64 $\times$ 768} 
= 245K params, materializing to 2.36M -- a 10$\times$ overhead. Used in Frobenius 
check inside \texttt{copy\_weights}.

\subsubsection{The `raise` Bug}

See Section 4.10.

\subsubsection{Softpick Masked Rows}

As noted, if a row is entirely masked, output is zero.

\subsubsection{The `\_kernel\_linear\_attention` Python Loop (Decoder)}

For $N = 128$, chunk\_size = 64: 2 iterations. Acceptable but not 
vectorized.

\subsubsection{`rff\_proj` is not persistent}

\begin{lstlisting}
self.register\_buffer("rff\_proj", torch.randn(...), persistent=False)
\end{lstlisting}

With \texttt{persistent=False}, \texttt{rff\_proj} is not saved in \texttt{state\_dict()}. On 
reload, it is re-initialized randomly. \textbf{This is a bug} -- the RFF 
projection must be preserved.

\subsection{Registry}

\begin{lstlisting}
@ATTENTION\_REGISTRY.register("orion")
@ATTENTION\_REGISTRY.register("orion\_attention")
@ATTENTION\_REGISTRY.register("orion\_module")
class ORIONAttention(BaseAttention):
    ...
\end{lstlisting}

\subsection{Compatibility with Recovery}

The \texttt{FactorisedLinear} supports LoRA adapters. The recovery pipeline:
\begin{enumerate}
\item Factorize base weights.
\item Add LoRA adapters.
\item Freeze base, train LoRA.
\end{enumerate}
\begin{lstlisting}

---
\end{lstlisting}

\begin{figure}[H]
\centering
\includegraphics[width=\linewidth]{figures/architecture.png}
\caption{High-level architectural diagram of orion.}
\label{fig:arch\_orion}
\end{figure}

